% Chương 6: Ma trận Kiểm chứng và Kịch bản Walkthrough Nghiệp vụ

Chương này cung cấp bộ công cụ kiểm chứng chất lượng và tính nhất quán toàn diện của đồ án. Bằng cách thiết lập Ma trận Truy vết Yêu cầu (\textit{Requirement Traceability Matrix - RTM}), Ma trận Quyền hạn Thao tác Dữ liệu (\textit{CRUD Matrix}), Ma trận Thực thi Quy tắc Nghiệp vụ (\textit{Rule-to-Enforcement Matrix}) và diễn giải 7 kịch bản nghiệp vụ thực tế (\textit{Scenario Walkthrough}), đồ án chứng minh rằng 100\% các yêu cầu và quy tắc nghiệp vụ đều được hiện thực hóa và bảo vệ toàn vẹn trong cơ sở dữ liệu.

% ----------------------------------------------------------------------
\section{Ma trận truy vết yêu cầu toàn diện (Requirement Traceability Matrix)}
\label{s:c6_rtm}

Ma trận RTM (Bảng~\ref{tab:c6_rtm}) liên kết xuyên suốt từ Yêu cầu chức năng / Quy tắc nghiệp vụ $\rightarrow$ Luồng nghiệp vụ $\rightarrow$ Phân hệ kiến trúc $\rightarrow$ Thực thể ý niệm $\rightarrow$ Bảng vật lý và Ràng buộc kiểm thử:

\begin{table}[!ht]
\centering
\caption{Ma trận truy vết yêu cầu toàn diện (RTM)}
\label{tab:c6_rtm}
\tiny
\begin{tabularx}{\textwidth}{p{2.0cm}p{1.1cm}p{1.6cm}p{2.4cm}p{2.8cm}>{\raggedright\arraybackslash}X}
\toprule
\textbf{Yêu cầu / BR} & \textbf{Luồng / UC} & \textbf{Phân hệ Module} & \textbf{Thực thể Logic} & \textbf{Bảng vật lý / Đối tượng} & \textbf{Kiểm thử / Ràng buộc} \\
\midrule
FR-007, FR-008, BR-P-002 & F02 / UC-02 & Registration & \texttt{Registration} & \texttt{participant.\allowbreak Registration} & \texttt{UQ (contest, participant)}, \texttt{TST-REG-001} \\
FR-010, FR-011, BR-P-005 & F03 / UC-03 & Film Asset & \texttt{FilmRoll}, \texttt{FilmFrame} & \texttt{film.FilmRoll}, \texttt{film.FilmFrame} & \texttt{UQ (roll, frame\_number)}, \texttt{TST-FRAME-001} \\
FR-013, FR-015, BR-P-006 & F04 / UC-04 & Submission & \texttt{Submission} & \texttt{submission.\allowbreak Submission}, \texttt{usp\_create\_\allowbreak submission} & \texttt{UQ (contest, frame\_id)}, \texttt{TST-SUB-001, 002} \\
FR-017, FR-018, BR-O-006 & F05 / UC-05 & Verification & \texttt{Verification\allowbreak Case}, \texttt{AIAnalysis\allowbreak Result} & \texttt{verification.\allowbreak VerificationCase}, \texttt{AIAnalysis\allowbreak Result} & \texttt{CK status}, \texttt{TST-VER-001} \\
FR-020, FR-021, BR-O-009 & F06 / UC-06 & Judging & \texttt{Judge\allowbreak Assignment}, \texttt{Evaluation} & \texttt{judging.\allowbreak JudgeAssignment}, \texttt{Evaluation} & \texttt{UQ (round, sub, judge)}, \texttt{TST-JDG-001, 002} \\
FR-025, FR-026, BR-O-010 & F07 / UC-07 & Result \& Award & \texttt{Result}, \texttt{Award\allowbreak Assignment} & \texttt{result.Result}, \texttt{usp\_finalize\_\allowbreak results} & \texttt{UQ rank}, \texttt{TST-RES-001} \\
FR-028, BR-O-011, BR-P-017 & F08 / UC-08 & Digital Archive & \texttt{ArchiveItem} & \texttt{archive.\allowbreak ArchiveItem} & \texttt{Snapshot integrity}, \texttt{TST-ARC-001} \\
FR-029 & All & Audit & \texttt{AuditLog} & \texttt{audit.AuditLog} & \texttt{Triggers}, \texttt{TST-AUD-001} \\
FR-030, BR-P-001 & All & Identity & \texttt{UserAccount}, \texttt{UserRole} & \texttt{iam.UserAccount}, \texttt{iam.UserRole} & \texttt{UQ email, username} \\
\bottomrule
\end{tabularx}
\end{table}

% ----------------------------------------------------------------------
\section{Ma trận quyền hạn thao tác dữ liệu (CRUD Matrix)}
\label{s:c6_crud}

Ma trận CRUD (Bảng~\ref{tab:c6_crud}) phân định quyền Tạo mới (C), Đọc (R), Cập nhật (U), Xóa (D) của từng phân hệ đối với các bảng CSDL:

\begin{table}[!ht]
\centering
\caption{Ma trận phân quyền thao tác dữ liệu (CRUD Matrix)}
\label{tab:c6_crud}
\scriptsize
\begin{tabularx}{\textwidth}{p{2.8cm}p{3.2cm}p{3.4cm}>{\raggedright\arraybackslash}X}
\toprule
\textbf{Phân hệ Module} & \textbf{Quyền Tạo mới (Create)} & \textbf{Quyền Đọc (Read)} & \textbf{Quyền Cập nhật (Update)} \\
\midrule
\textbf{Identity \& Access} & \texttt{UserAccount}, \texttt{UserRole} & \texttt{UserAccount}, \texttt{Role}, \texttt{UserRole} & \texttt{UserAccount}, \texttt{UserRole} \\
\textbf{Contest Management} & \texttt{Contest}, \texttt{Category}, \texttt{Round}, \texttt{Criterion}, \texttt{Award} & Các bảng cấu hình cuộc thi & Các bảng cấu hình cuộc thi \\
\textbf{Registration} & \texttt{Registration} & \texttt{Contest}, \texttt{ParticipantProfile} & \texttt{Registration} \\
\textbf{Film Asset Management} & \texttt{FilmRoll}, \texttt{FilmFrame} & Danh mục máy ảnh, ống kính, lab & \texttt{FilmRoll}, \texttt{FilmFrame} \\
\textbf{Submission Management} & \texttt{Submission} & \texttt{Registration}, \texttt{FilmFrame} & \texttt{Submission} (trạng thái) \\
\textbf{Verification} & \texttt{VerificationCase}, \texttt{AIAnalysisResult} & \texttt{Submission}, \texttt{AIAnalysisResult} & \texttt{VerificationCase} \\
\textbf{Judging} & \texttt{JudgeAssignment}, \texttt{Evaluation} & Bài thi, vòng chấm, tiêu chí & \texttt{JudgeAssignment}, \texttt{Evaluation} \\
\textbf{Result \& Award} & \texttt{Result}, \texttt{AwardAssignment} & \texttt{Evaluation}, \texttt{AwardDefinition} & \texttt{Result} (trạng thái) \\
\textbf{Digital Archive} & \texttt{ArchiveItem} & \texttt{Result}, \texttt{Submission}, \texttt{FilmFrame} & Không sửa đổi (Bất biến) \\
\textbf{Audit \& Reporting} & \texttt{AuditLog} & Toàn bộ các bảng CSDL & Không cập nhật \\
\bottomrule
\end{tabularx}
\end{table}

% ----------------------------------------------------------------------
\section{Kiểm chứng 7 kịch bản nghiệp vụ thực tế (Scenario Walkthrough)}
\label{s:c6_scenarios}

Hệ thống thiết kế đã được kiểm chứng thành công qua 7 kịch bản nghiệp vụ phức tạp:

\begin{enumerate}
    \item \textbf{Kịch bản S1 -- Nộp bài Hợp lệ (Normal Submission)}: Thí sinh nộp khung hình số 12 của cuộn Kodak Portra 400 vào cuộc thi trước hạn chót. Thủ tục \texttt{usp\_create\_submission} kiểm tra hợp lệ và khởi tạo bản ghi \texttt{Submission} ở trạng thái \texttt{PendingVerification}.
    \item \textbf{Kịch bản S2 -- Xử lý Cảnh báo AI Trùng lặp (Duplicate / AI Flag Review)}: Module AI phân tích ảnh và ghi nhận cờ \texttt{FLAGGED\_SUSPICIOUS} kèm độ tin cậy $0.9450$ vào bảng \texttt{AIAnalysisResult}. Ban tổ chức mở hàng đợi kiểm duyệt, đối soát thủ công và đưa ra quyết định từ chối (\texttt{Rejected}) qua thủ tục \texttt{usp\_record\_verification\_decision}.
    \item \textbf{Kịch bản S3 -- Chặn Nộp bài Quá hạn hoặc Sai lệch Xuất xứ (Late/Invalid Submission)}: Thí sinh cố tình nộp bài khi thời gian nộp bài đã đóng (\texttt{submission\_close\_at}), hoặc nộp khung hình không thuộc quyền sở hữu của mình. Thủ tục \texttt{usp\_create\_submission} tự động bắt lỗi và hủy giao dịch (\texttt{ROLLBACK}).
    \item \textbf{Kịch bản S4 -- Chấm thi Đa vòng Độc lập (Multi-round Judging)}: Bài thi vượt qua Vòng Sơ khảo được đưa vào Vòng Chung khảo. Điểm số của Vòng 1 và Vòng 2 được lưu trữ tại các bản ghi \texttt{Evaluation} riêng biệt theo từng \texttt{round\_id}, hoàn toàn không bị ghi đè.
    \item \textbf{Kịch bản S5 -- Đảm bảo Tính Duy nhất của Phiếu Chấm (Judge Uniqueness)}: Một giám khảo cố tình nộp phiếu chấm thứ hai cho cùng một bài thi trong cùng một vòng. Ràng buộc \texttt{UQ\_judging\_Evaluation\_round\_sub\_judge} lập tức kích hoạt lỗi vi phạm khóa duy nhất.
    \item \textbf{Kịch bản S6 -- Khóa Kết quả và Trao giải (Result Finalization)}: Sau khi kết thúc thời gian chấm thi, thủ tục \texttt{usp\_finalize\_results\_for\_round} tự động tính tổng điểm trung bình có trọng số, xếp hạng phân biệt từ 1 đến $N$, gán giải thưởng và khóa kết quả chuyển sang trạng thái \texttt{Published}.
    \item \textbf{Kịch bản S7 -- Tự động Lưu trữ Di sản Số (Digital Archiving)}: Ngay sau khi công bố giải, tác phẩm đạt giải Nhất được chụp toàn bộ siêu dữ liệu (cuộn phim, thông số scan, nhận xét của giám khảo, điểm số) lưu vào bảng \texttt{ArchiveItem}, sẵn sàng phục vụ triển lãm trực tuyến lâu dài.
\end{enumerate}

% ----------------------------------------------------------------------
\section{Kết quả thực thi toàn diện bộ 35 kịch bản kiểm thử T-SQL tự động}
\label{s:c6_test_results}

Toàn bộ 35 kịch bản kiểm thử T-SQL tự động trong thư mục \texttt{tests/} đã được thực thi và xác nhận vượt qua 100\% trên môi trường Microsoft SQL Server 2022 CU24 containerized (\texttt{tkcsdl-sqlserver}, cổng 14333) kết nối trực tiếp qua SSMS 22. Bộ test được chia thành 4 nhóm kiểm thử chuyên sâu (Bảng~\ref{tab:c6_tests}):

\begin{table}[!ht]
\centering
\caption{Kết quả thực thi bộ 35 kịch bản kiểm thử tự động T-SQL trên SQL Server 2022}
\label{tab:c6_tests}
\tiny
\begin{tabularx}{\textwidth}{>{\bfseries\ttfamily}p{2.6cm}>{\raggedright\arraybackslash}X>{\raggedright\arraybackslash}p{2.6cm}}
\toprule
Mã Kịch bản & Nội dung kiểm thử nghiệp vụ & Kết quả thực thi \\
\midrule
\multicolumn{3}{l}{\textbf{Nhóm 1: Ràng buộc Toàn vẹn, Khóa Nghiệp vụ và Miền giá trị}} \\
\midrule
TST-REG-001 & Chặn đăng ký trùng lặp cho cùng một thí sinh trong cuộc thi & \textbf{PASS} (Bắt lỗi UQ) \\
TST-FRAME-001 & Chặn trùng lặp số frame trong cùng một cuộn phim & \textbf{PASS} (Bắt lỗi UQ) \\
TST-UQ-001 & Kiểm chứng tính duy nhất của toàn bộ khóa nghiệp vụ tự nhiên & \textbf{PASS} (Toàn vẹn UQ) \\
TST-INT-001 & Chặn vi phạm khóa ngoại không tồn tại giữa các phân hệ & \textbf{PASS} (Bắt lỗi FK) \\
TST-FK-001 & Kiểm chứng liên kết khóa ngoại và toàn vẹn tham chiếu & \textbf{PASS} (Hợp lệ 100\%) \\
TST-INT-002 & Chặn trạng thái thực thể không thuộc tập giá trị cho phép & \textbf{PASS} (Bắt lỗi Check) \\
TST-INT-003 & Chặn trọng số tiêu chí chấm sai lệch hoặc vượt quá 100\% & \textbf{PASS} (Bắt lỗi Check) \\
TST-INT-004 & Chặn cửa sổ thời gian nộp bài bị đảo ngược & \textbf{PASS} (Bắt lỗi Check) \\
TST-BOUNDARY-001 & Kiểm thử giá trị biên trạng thái, thang điểm và ngày tháng & \textbf{PASS} (Kiểm soát biên) \\
TST-RET-001 & Ngăn chặn xóa vật lý dữ liệu lịch sử và kết quả đã lưu trữ & \textbf{PASS} (Restricted Delete) \\
TST-SEED-001 & Kiểm chứng tính nhất quán tập dữ liệu tham chiếu nền tảng & \textbf{PASS} (Dữ liệu nền tảng) \\
TST-SEED-002 & Kiểm chứng tính toàn vẹn của đồ thị dữ liệu mẫu & \textbf{PASS} (Đồ thị nhất quán) \\
\midrule
\multicolumn{3}{l}{\textbf{Nhóm 2: Thủ tục Lưu trữ Nghiệp vụ và Kiểm soát Giao dịch ACID}} \\
\midrule
TST-SUB-001 & Nộp bài dự thi hợp lệ thành công qua thủ tục lưu trữ & \textbf{PASS} (Tạo bản ghi) \\
TST-SUB-002 & Chặn nộp bài quá hạn chót hoặc không sở hữu khung hình & \textbf{PASS} (Bắt lỗi SP) \\
TST-SUB-003 & Ma trận kiểm thử toàn diện các điều kiện biên của thủ tục nộp bài & \textbf{PASS} (Xử lý đa kịch bản) \\
TST-TXN-001 & Tự động hoàn nguyên giao dịch (Rollback) khi phát sinh lỗi & \textbf{PASS} (Bảo toàn ACID) \\
TST-VER-001 & Thẩm định bài thi có cảnh báo AI và cập nhật trạng thái & \textbf{PASS} (Duyệt con người) \\
TST-VER-002 & Cổng thẩm định con người và chặn đổi trạng thái kết thúc & \textbf{PASS} (Bảo vệ Terminal) \\
TST-JDG-001 & Chặn giám khảo nộp phiếu chấm lần hai trong cùng một vòng & \textbf{PASS} (Bắt lỗi UQ) \\
TST-JDG-002 & Chấm thi đa vòng với tiêu chí và trọng số độc lập & \textbf{PASS} (Độc lập dữ liệu) \\
TST-FN-001 & Hàm vô hướng \texttt{fn\_round\_max\_total\_score} tính điểm trần & \textbf{PASS} (Khớp tổng điểm) \\
TST-RES-001 & Tự động tổng hợp điểm, xếp hạng phân biệt và khóa kết quả & \textbf{PASS} (Xếp hạng chuẩn) \\
TST-RES-002 & Kiểm chứng tính lũy đẳng (Idempotency) khi chạy lại chốt kết quả & \textbf{PASS} (Kết quả ổn định) \\
TST-ARC-001 & Kiểm chứng tính bất biến của Snapshot lưu trữ khi Profile đổi & \textbf{PASS} (Snapshot chuẩn) \\
TST-AUD-001 & Trigger tự động ghi nhật ký kiểm toán khi đổi trạng thái bài thi & \textbf{PASS} (Ghi log chuẩn) \\
\midrule
\multicolumn{3}{l}{\textbf{Nhóm 3: Tính năng Nâng cao (Template, Siêu dữ liệu số, Immutability, Trao giải)}} \\
\midrule
TST-TMPL-001 & Nhân bản hoàn chỉnh cấu hình cuộc thi từ Contest Template chuẩn & \textbf{PASS} (Cloning Hierarchy) \\
TST-ASSET-001 & Cấu trúc siêu dữ liệu số, mã băm SHA-256 (64 hex) và bit depth & \textbf{PASS} (Toàn vẹn Metadata) \\
TST-ARC-002 & Chặn tạo lưu trữ từ kết quả Draft/chưa Finalized và chặn trùng & \textbf{PASS} (Bắt lỗi 55001/2) \\
TST-IMMUT-001 & Chặn sửa/xóa trực tiếp Result đã Finalized và ArchiveItem qua Trigger & \textbf{PASS} (Bắt lỗi 54011/55010) \\
TST-PUB-001 & Kiểm tra tính đầy đủ cấu hình trước khi Publish cuộc thi & \textbf{PASS} (Bắt lỗi 51202-5) \\
TST-AWD-001 & Chặn trao giải thưởng sai lệch danh mục qua Trigger và Procedure & \textbf{PASS} (Bắt lỗi 54103/5) \\
TST-VIEW-001 & Kiểm tra tính toàn vẹn runtime và hợp lệ cú pháp JSON của các View & \textbf{PASS} (5 Views hợp lệ) \\
\midrule
\multicolumn{3}{l}{\textbf{Nhóm 4: Mô phỏng Tranh chấp Đồng thời Hai Phiên (Concurrency Race Conditions)}} \\
\midrule
TST-CONC-001 & Tranh chấp đồng thời khi 2 phiên cùng đăng ký một cuộc thi & \textbf{PASS} (Chặn trùng UQ) \\
TST-CONC-002 & Tranh chấp đồng thời khi 2 phiên cùng nộp một khung hình phim & \textbf{PASS} (Chặn nộp đúp) \\
TST-CONC-003 & Tranh chấp đồng thời khi giám khảo gửi 2 phiếu chấm cùng lúc & \textbf{PASS} (Bảo toàn duy nhất) \\
\bottomrule
\end{tabularx}
\end{table}

Kết quả thực nghiệm trên container SQL Server 2022 CU24 xác nhận rằng hệ thống cơ sở dữ liệu đã đạt mức hoàn thiện 100\%: tất cả các quy tắc nghiệp vụ từ cơ bản đến phức tạp (tính lũy đẳng, cấu hình mẫu tái sử dụng, bảo vệ bất biến dữ liệu di sản số, phòng chống tranh chấp đồng thời) đều được thực thi chặt chẽ, chứng minh cơ sở dữ liệu hoàn toàn sẵn sàng cho môi trường thực tế.
